\documentclass[aspectratio=169,11pt]{beamer}
\usepackage{../../../shared/theme/esmad_beamer_theme}
\usepackage{amsmath,amssymb}
\usepackage{booktabs}

\title[Counterfactual Time Series]{Counterfactual Time-Series Methods}
\subtitle{Interrupted Time Series and Synthetic Control}
\date{\today}

\newcommand{\E}{\mathbb{E}}

\begin{document}

\begin{frame}
\titlepage
\end{frame}

\begin{frame}{Learning Goals}
\begin{itemize}
\item distinguish descriptive structural breaks from causal intervention effects;
\item formulate interrupted time-series models with level and slope changes;
\item account for seasonality, serial correlation, and concurrent shocks;
\item construct synthetic controls from donor units and assess pre-treatment fit;
\item use placebo units and placebo intervention dates as falsification devices;
\item compare when interrupted time series, controlled ITS, and synthetic control are credible.
\end{itemize}
\end{frame}

\begin{frame}{Outline}
\tableofcontents
\end{frame}

\section{Counterfactual Time Series}

\begin{frame}{The Core Causal Question}
Suppose an intervention occurs at time \(T_0\).

For \(t>T_0\), we observe
\[
Y_t(1),
\]
but we need the missing untreated trajectory
\[
Y_t(0).
\]

The causal effect at time \(t\) is
\[
\tau_t=Y_t(1)-Y_t(0).
\]

\begin{alertblock}{The challenge}
Time itself does not provide the counterfactual. We must model or construct what would have happened without intervention.
\end{alertblock}
\end{frame}

\begin{frame}{Before--After Is Usually Not Enough}
A simple comparison
\[
\bar Y_{post}-\bar Y_{pre}
\]
attributes every temporal change to treatment.

That can confound treatment with:
\begin{itemize}
\item secular trends;
\item seasonality;
\item macroeconomic shocks;
\item measurement changes;
\item unrelated structural breaks.
\end{itemize}
\end{frame}

\section{Interrupted Time Series}

\begin{frame}{Segmented Regression}
A basic interrupted time-series model is
\[
Y_t
=
\beta_0
+
\beta_1 Time_t
+
\beta_2 Intervention_t
+
\beta_3 PostTime_t
+
\varepsilon_t.
\]

Where:
\begin{itemize}
\item \(\beta_1\): pre-intervention slope;
\item \(\beta_2\): immediate level change;
\item \(\beta_3\): slope change after intervention.
\end{itemize}
\end{frame}

\begin{frame}{Level Change vs Slope Change}
An intervention may cause:
\begin{itemize}
\item an immediate jump with no trend change;
\item no immediate jump but a gradual trend change;
\item both;
\item a transient effect followed by reversion.
\end{itemize}

\begin{block}{Interpretation}
The post-treatment counterfactual is an extrapolation of the untreated data-generating process, not simply the pre-period mean.
\end{block}
\end{frame}

\begin{frame}{Key ITS Identification Assumption}
The pre-intervention process must provide a credible basis for the untreated post-intervention trajectory.

This requires ruling out alternative explanations for the break:
\begin{itemize}
\item no concurrent policy shock;
\item no abrupt measurement change;
\item no composition change;
\item no anticipatory response contaminating the pre-period.
\end{itemize}
\end{frame}

\section{Time-Series Dependence}

\begin{frame}{Autocorrelation Matters}
Residuals often satisfy
\[
\Cov(\varepsilon_t,\varepsilon_{t-k})\neq 0.
\]

Ignoring serial correlation can understate uncertainty around the intervention effect.

Possible responses include:
\begin{itemize}
\item explicit AR error structures;
\item HAC-style standard errors;
\item state-space formulations;
\item block bootstrap methods.
\end{itemize}
\end{frame}

\begin{frame}{Seasonality}
Suppose outcomes vary systematically by month or quarter.

A segmented regression may include seasonal terms:
\[
Y_t
=
\beta_0+\beta_1 Time_t+\beta_2 D_t+\beta_3 PostTime_t
+
\sum_{m=2}^{12}\delta_m Month_{mt}
+
\varepsilon_t.
\]

\begin{alertblock}{Threat}
If intervention timing coincides with seasonal peaks or troughs, naive before--after comparisons can be badly misleading.
\end{alertblock}
\end{frame}

\begin{frame}{Structural Breaks vs Treatment Effects}
A detected break at \(T_0\) shows that the data-generating process changed.

It does not, by itself, show that treatment caused the break.

\begin{block}{Causal requirement}
The intervention must be the most credible explanation for the deviation from the untreated counterfactual.
\end{block}
\end{frame}

\section{Controlled Interrupted Time Series}

\begin{frame}{Adding a Comparison Series}
Suppose a comparable untreated series \(C_t\) is available.

A controlled ITS uses that series to absorb common shocks:
\[
Y_t
=
\alpha
+
\beta_1 Time_t
+
\beta_2 TreatSeries
+
\beta_3 Post_t
+
\beta_4(TreatSeries\times Post_t)
+
\cdots
\]

This creates a time-series analogue of comparative counterfactual reasoning.
\end{frame}

\begin{frame}{When a Control Series Helps}
A control series can improve credibility when it:
\begin{itemize}
\item shares common shocks with the treated series;
\item is not itself affected by treatment;
\item has similar measurement quality;
\item follows a stable relationship with the treated series pre-treatment.
\end{itemize}

\begin{alertblock}{Spillovers}
If treatment affects the control series, the counterfactual is contaminated.
\end{alertblock}
\end{frame}

\section{Synthetic Control}

\begin{frame}{Why Synthetic Control?}
Sometimes there is one treated unit and many potential comparison units.

No single donor may be sufficiently similar.

Synthetic control constructs
\[
\widehat Y_{1t}(0)
=
\sum_{j=2}^{J+1}w_jY_{jt}
\]
with
\[
w_j\geq 0,
\qquad
\sum_jw_j=1.
\]

The weights are chosen to reproduce pre-treatment characteristics and outcomes.
\end{frame}

\begin{frame}{The Donor Pool}
The donor pool should contain units that:
\begin{itemize}
\item were not treated during the relevant period;
\item were not substantially affected by spillovers;
\item share plausible structural relationships with the treated unit;
\item have reliable pre-treatment data.
\end{itemize}

\begin{alertblock}{Design decision}
Donor selection is part of identification. A large donor pool is not automatically a good donor pool.
\end{alertblock}
\end{frame}

\begin{frame}{Pre-Treatment Fit}
Synthetic control credibility depends heavily on how well the weighted donors reproduce the treated unit before intervention.

A common diagnostic is pre-treatment RMSPE:
\[
RMSPE_{pre}
=
\sqrt{
\frac{1}{T_0}
\sum_{t<T_0}
(Y_{1t}-\widehat Y_{1t}(0))^2
}.
\]

Poor pre-fit weakens the post-treatment counterfactual.
\end{frame}

\begin{frame}{Post-Treatment Effect Path}
After treatment,
\[
\widehat\tau_t
=
Y_{1t}
-
\widehat Y_{1t}(0).
\]

Inspect:
\begin{itemize}
\item onset timing;
\item persistence;
\item cumulative effect;
\item whether the gap is large relative to pre-treatment mismatch.
\end{itemize}
\end{frame}

\section{Placebo Inference}

\begin{frame}{Placebo Units}
Pretend each donor unit was treated at the same date and construct a synthetic control for it.

Then compare the treated unit's post-treatment gap with placebo gaps.

\begin{block}{Purpose}
This asks whether an effect of the observed magnitude is unusual relative to units that were not treated.
\end{block}
\end{frame}

\begin{frame}{Placebo Intervention Dates}
Move the treatment date into the pre-treatment period.

A large apparent effect before actual intervention can indicate:
\begin{itemize}
\item unstable pre-treatment relationships;
\item overfitting;
\item structural breaks unrelated to treatment;
\item model misspecification.
\end{itemize}
\end{frame}

\begin{frame}{Post/Pre Fit Ratios}
One descriptive diagnostic compares post-treatment and pre-treatment fit:
\[
\frac{RMSPE_{post}}{RMSPE_{pre}}.
\]

A treated unit with a much larger ratio than placebo units can be suggestive.

\begin{alertblock}{Caution}
These placebo distributions are design-based diagnostics, not automatic classical p-values.
\end{alertblock}
\end{frame}

\section{Failure Modes}

\begin{frame}{Concurrent Shocks}
Suppose treatment begins during:
\begin{itemize}
\item a recession;
\item a supply-chain disruption;
\item another policy change;
\item a measurement-system migration.
\end{itemize}

If these shocks affect the treated unit differently, the estimated intervention effect may absorb them.
\end{frame}

\begin{frame}{Donor Contamination}
A donor unit is problematic if it:
\begin{itemize}
\item later adopts the treatment;
\item receives spillovers;
\item experiences a unique structural break;
\item changes measurement definitions.
\end{itemize}

\begin{block}{Practical rule}
Audit donor histories, not just statistical fit.
\end{block}
\end{frame}

\begin{frame}{Overfitting the Pre-Period}
A synthetic control can match pre-treatment data extremely closely yet extrapolate poorly after intervention.

Possible reasons:
\begin{itemize}
\item too many predictors relative to pre-period length;
\item unstable donor relationships;
\item accidental fit to noise;
\item structural change unrelated to treatment.
\end{itemize}
\end{frame}

\section{Choosing Between Designs}

\begin{frame}{ITS vs Synthetic Control}
\begin{tabular}{p{3.3cm}p{4.4cm}p{4.4cm}}
\toprule
& \textbf{Interrupted time series} & \textbf{Synthetic control} \\
\midrule
Counterfactual source & Extrapolated own history & Weighted donor units \\
Main strength & Rich temporal structure & External comparison trajectory \\
Main risk & Concurrent structural break & Poor donor support \\
Best setting & Long stable pre-period & One/few treated units, good donors \\
\bottomrule
\end{tabular}
\end{frame}

\begin{frame}{When to Prefer Controlled ITS}
Controlled ITS is attractive when:
\begin{itemize}
\item a natural comparison series exists;
\item common shocks are important;
\item treatment affects only one series;
\item both series have long pre- and post-periods.
\end{itemize}

It can be simpler and more interpretable than a synthetic donor construction when a credible control is obvious.
\end{frame}

\section{Applied Workflow}

\begin{frame}{Example: Regional Pricing Policy}
Suppose one region changes pricing policy in January 2026.

Questions:
\begin{itemize}
\item Was there already a demand trend before January?
\item Is January seasonally unusual?
\item Did competitors change pricing simultaneously?
\item Are neighboring regions unaffected and suitable donors?
\item Does the synthetic trajectory fit well before treatment?
\item Are post-treatment gaps larger than placebo gaps?
\end{itemize}
\end{frame}

\begin{frame}{A Credible Counterfactual Time-Series Workflow}
\begin{enumerate}
\item Define intervention date and estimand.
\item Plot raw series and inspect seasonality.
\item Specify the untreated time-series process.
\item Audit concurrent shocks and anticipation.
\item Add control series or donor units where possible.
\item Diagnose pre-treatment fit.
\item Model serial correlation.
\item Run placebo units and placebo dates.
\item Report effect paths and sensitivity to donor/model choices.
\end{enumerate}
\end{frame}

\begin{frame}{Synthesis: Building a Missing Trajectory}
\begin{itemize}
\item Causal time-series analysis is fundamentally a counterfactual forecasting problem.
\item Interrupted time series extrapolates the treated unit's own untreated process.
\item Controlled ITS strengthens the design with a comparison series.
\item Synthetic control constructs the untreated trajectory from weighted donors.
\item Pre-treatment fit is necessary but not sufficient for credibility.
\item Serial correlation, seasonality, spillovers, and concurrent shocks must be addressed explicitly.
\item Placebo units and placebo dates are central falsification tools.
\end{itemize}
\end{frame}

\begin{frame}{Selected References}
\small
\begin{itemize}
\item Bernal, J. L., Cummins, S., and Gasparrini, A. (2017). Interrupted Time Series Regression for the Evaluation of Public Health Interventions. \textit{International Journal of Epidemiology}, 46(1), 348--355.
\item Abadie, A., Diamond, A., and Hainmueller, J. (2010). Synthetic Control Methods for Comparative Case Studies. \textit{Journal of the American Statistical Association}, 105(490), 493--505.
\item Abadie, A., Diamond, A., and Hainmueller, J. (2015). Comparative Politics and the Synthetic Control Method. \textit{American Journal of Political Science}, 59(2), 495--510.
\item Abadie, A. (2021). Using Synthetic Controls: Feasibility, Data Requirements, and Methodological Aspects. \textit{Journal of Economic Literature}, 59(2), 391--425.
\item Box, G. E. P., Jenkins, G. M., Reinsel, G. C., and Ljung, G. M. (2015). \textit{Time Series Analysis: Forecasting and Control}, 5th ed. Wiley.
\end{itemize}
\end{frame}

\contactslide

\end{document}
